% Plan v5 / item F: the opening effect, migrated from the main text.
\section{Where the uncertified interval comes from}\label{app:opening}

% analysis/opening_effect.py -> results/opening_effect{,_tinycomma18b,_pleias3b,_kl3m37b}{,_token}.csv
Measured on the 758 reference passages under four safe models from three openly licensed corpora,
$k_{\mathrm{crit}}/s$ has a median of $4.25$ to $6.10$. That figure should not be quoted alone. The
running maximum in Eq.~\eqref{eq:rates} is attained at the \emph{first} token of the work in
$87.7\%$ to $90.5\%$ of cases and within the opening tenth in $96.7\%$ to $97.5\%$, because the
bucket must pay for the opening out of a single step's allowance. An adversary who supplies one
token of genuine prefix removes it: the same ratio falls to $1.32$--$1.48$ and stays there when
$5$, $10$ or $20$ tokens are supplied. The uncertified protection is real, and it is almost
entirely a property of where a generation starts.

The choice of denominator flatters this effect without manufacturing it, and the difference is
worth stating rather than asserting robustness. Repeating every measurement per token leaves the
ratios after a supplied prefix unchanged --- $1.35$--$1.41$ at one token and $1.30$--$1.45$ at ten,
against $1.32$--$1.48$ and $1.30$--$1.47$ per character --- but lowers the whole-work median to
$3.88$--$5.98$ and the fraction binding at the very first token to $82.6\%$--$87.7\%$, with
$93.0\%$--$95.4\%$ inside the opening tenth. The opening effect is smaller than the per-character
figures alone suggest and remains the dominant term under either denominator.

This also explains a gap the audit measured without accounting for it. Oracle-prefix windows
recover $0.2301$ of a passage at $k=5$ against $0.0925$ for a single query in an earlier audit of
this mechanism \citep{vijayavallabh2026audit}, and $0.097$ against $0.012$ at $k=3$ here
(Table~\ref{tab:orders}); the oracle strategy is precisely one that supplies genuine prefixes,
and Proposition~\ref{prop:outrun} says that is what dismantles the opening effect. It unifies the
prefix debt with the theorem as well: $\delta$ raises $k_{\mathrm{crit}}$ and the running maximum
binds at the opening, so the untheorised heuristic and the theorem's binding constraint are the
same phenomenon seen twice, and an adversary who owns the prefix defeats both.


\begin{figure}[h]
\centering
\includegraphics[width=0.6\textwidth]{opening_effect.pdf}
\caption{The gap between certifying and protecting, $k_{\mathrm{crit}}/s(x)$, against the number of
genuine prefix tokens the adversary supplies. It is large only for a generation that starts from
nothing.}
\label{fig:opening}
\end{figure}
